%==============================================================================
%  lesson03.tex
%  The mixed formulation, saddle point structure, and the inf-sup condition
%  Repository: micropolar-fem
%  Author: M. Javed Bukhari
%
%  Requires macros.tex in the same folder.
%  Compile with: pdflatex lesson03.tex   (run it twice, for the contents page)
%==============================================================================

\documentclass[11pt,a4paper]{article}

% Shared macros for the theory notes.
% Include with % Shared macros for the theory notes.
% Include with % Shared macros for the theory notes.
% Include with \input{macros} after \begin{document} preamble packages.

\newcommand{\Om}{\Omega}
\newcommand{\dOm}{\partial\Omega}
\newcommand{\dx}{\,\mathrm{d}x}
\newcommand{\ds}{\,\mathrm{d}s}

\newcommand{\Ltwo}{L^2(\Om)}
\newcommand{\Ltwozero}{L^2_0(\Om)}
\newcommand{\Hone}{H^1(\Om)}
\newcommand{\Honezero}{H^1_0(\Om)}

\newcommand{\normL}[1]{\left\lVert #1 \right\rVert_{L^2(\Om)}}
\newcommand{\normH}[1]{\left\lVert #1 \right\rVert_{H^1(\Om)}}
\newcommand{\ip}[2]{\left( #1 , #2 \right)}

\newcommand{\grad}{\nabla}
\newcommand{\lap}{\Delta}
\newcommand{\dn}[1]{\frac{\partial #1}{\partial n}}

\newtheorem{theorem}{Theorem}[section]
\newtheorem{lemma}[theorem]{Lemma}
\newtheorem{proposition}[theorem]{Proposition}
\theoremstyle{definition}
\newtheorem{definition}[theorem]{Definition}
\newtheorem{remark}[theorem]{Remark}
 after \begin{document} preamble packages.

\newcommand{\Om}{\Omega}
\newcommand{\dOm}{\partial\Omega}
\newcommand{\dx}{\,\mathrm{d}x}
\newcommand{\ds}{\,\mathrm{d}s}

\newcommand{\Ltwo}{L^2(\Om)}
\newcommand{\Ltwozero}{L^2_0(\Om)}
\newcommand{\Hone}{H^1(\Om)}
\newcommand{\Honezero}{H^1_0(\Om)}

\newcommand{\normL}[1]{\left\lVert #1 \right\rVert_{L^2(\Om)}}
\newcommand{\normH}[1]{\left\lVert #1 \right\rVert_{H^1(\Om)}}
\newcommand{\ip}[2]{\left( #1 , #2 \right)}

\newcommand{\grad}{\nabla}
\newcommand{\lap}{\Delta}
\newcommand{\dn}[1]{\frac{\partial #1}{\partial n}}

\newtheorem{theorem}{Theorem}[section]
\newtheorem{lemma}[theorem]{Lemma}
\newtheorem{proposition}[theorem]{Proposition}
\theoremstyle{definition}
\newtheorem{definition}[theorem]{Definition}
\newtheorem{remark}[theorem]{Remark}
 after \begin{document} preamble packages.

\newcommand{\Om}{\Omega}
\newcommand{\dOm}{\partial\Omega}
\newcommand{\dx}{\,\mathrm{d}x}
\newcommand{\ds}{\,\mathrm{d}s}

\newcommand{\Ltwo}{L^2(\Om)}
\newcommand{\Ltwozero}{L^2_0(\Om)}
\newcommand{\Hone}{H^1(\Om)}
\newcommand{\Honezero}{H^1_0(\Om)}

\newcommand{\normL}[1]{\left\lVert #1 \right\rVert_{L^2(\Om)}}
\newcommand{\normH}[1]{\left\lVert #1 \right\rVert_{H^1(\Om)}}
\newcommand{\ip}[2]{\left( #1 , #2 \right)}

\newcommand{\grad}{\nabla}
\newcommand{\lap}{\Delta}
\newcommand{\dn}[1]{\frac{\partial #1}{\partial n}}

\newtheorem{theorem}{Theorem}[section]
\newtheorem{lemma}[theorem]{Lemma}
\newtheorem{proposition}[theorem]{Proposition}
\theoremstyle{definition}
\newtheorem{definition}[theorem]{Definition}
\newtheorem{remark}[theorem]{Remark}


\newcommand{\Vdiv}{\mathbf{V}_{\mathrm{div}}}
\newcommand{\Vsp}{\mathbf{V}}
\newcommand{\Qsp}{Q}
\newcommand{\Wsp}{W}
\newcommand{\Xsp}{X}
\newcommand{\Zsp}{Z}
\newcommand{\bff}{\mathbf{f}}
\newcommand{\bzero}{\mathbf{0}}

\title{The Mixed Formulation and the Inf-Sup Condition:\\
Why Coercivity Is Not Enough}
\author{M. Javed Bukhari}
\date{\today}

\begin{document}
\maketitle

\begin{abstract}
In Lesson 2 I proved that the linear stationary micropolar system is well posed
on the space of divergence free velocity fields. That result is correct but it
cannot be used for computation, because there is no convenient finite element
basis consisting of exactly divergence free functions. These notes therefore go
back to the full problem, keep the pressure as an unknown, and impose
incompressibility weakly. The price is that the resulting bilinear form is no
longer coercive, so Lax--Milgram no longer applies. It has to be replaced by
Brezzi's theorem, which asks for two conditions instead of one: coercivity on
the kernel, which is exactly what Lesson 2 already proved, and a new condition
called the inf-sup condition. I set up all the spaces carefully, explain what
every term in the inf-sup condition means and why it has the shape it does, and
then work through five tasks establishing the structure of the mixed problem.
\end{abstract}

\tableofcontents

%==============================================================================
\section{Why we are changing the formulation}
%==============================================================================

Lesson 2 worked on
\[
  \Vdiv = \left\{ \bv \in [\Honezero]^{2} : \dive \bv = 0 \right\} ,
\]
and the pressure disappeared from the weak form because the term
$-\int_\Om p \,\dive \bv \dx$ vanishes identically when $\bv$ is divergence
free. That gave a clean coercive problem and Lax--Milgram applied directly.

The difficulty is entirely computational. To discretise a problem posed on
$\Vdiv$ we would need a finite element basis whose functions are
\emph{exactly} divergence free, that is, piecewise polynomial vector fields
$\bv_{h}$ with $\dive \bv_{h} = 0$ pointwise. Such bases do exist, the
Scott--Vogelius element being the best known example, but they require special
meshes or high polynomial degree, and they are not what standard finite element
software provides. If we take the ordinary vector valued $P_{2}$ basis, not a
single one of its functions is divergence free.

So we go back to the full system, keep the pressure $p$ as a genuine unknown,
and impose $\dive \bu = 0$ \emph{weakly}, meaning as an equation tested against
pressure functions rather than as a property built into the space. This is
called the mixed formulation, and it is what every practical finite element code
for incompressible flow actually solves.

The price is the subject of these notes. Once the space is enlarged, the
bilinear form stops being coercive, and the theory has to change.

%==============================================================================
\section{The spaces}
%==============================================================================

Since the whole lesson is about which space each unknown lives in, I want to set
all of them out carefully, with definitions, before anything else happens.

\subsection{Recalling $L^{2}$ and $H^{1}$}

\begin{definition}[The space $\Ltwo$]
$\Ltwo$ is the set of measurable functions $\phi : \Om \to \R$ with
\[
  \int_\Om \abs{\phi}^{2} \dx < \infty ,
\]
where two functions are regarded as the same element if they agree outside a set
of measure zero. It is equipped with the inner product and norm
\[
  \ip{\phi}{\psi} = \int_\Om \phi \psi \dx ,
  \qquad
  \normL{\phi} = \ip{\phi}{\phi}^{1/2} .
\]
\end{definition}

The intuition, from Lesson 1, is that $\normL{\phi}^{2}$ is a total squared
magnitude. When $\phi$ is a velocity it is kinetic energy up to a constant, so
membership of $\Ltwo$ is the statement that the field has finite energy.

The identification of functions agreeing almost everywhere is needed for
$\normL{\cdot}$ to be a genuine norm rather than a seminorm. If we did not
identify them, a function which is zero except at one point would have zero norm
without being the zero function, and the axiom
$\normL{\phi} = 0 \Rightarrow \phi = 0$ would fail.

\begin{definition}[The spaces $\Hone$ and $\Honezero$]
$\Hone$ is the set of $\phi \in \Ltwo$ whose weak first derivatives also lie in
$\Ltwo$, with norm
\[
  \normH{\phi}^{2} = \normL{\phi}^{2} + \normL{\grad \phi}^{2} .
\]
The subspace $\Honezero$ is the closure of $\Cinfc$ in this norm, which for a
Lipschitz domain coincides with the set of $\Hone$ functions whose trace on
$\dOm$ vanishes.
\end{definition}

Physically, $\normL{\grad \phi}^{2}$ measures how fast the field varies, and for
a velocity field it is proportional to the rate of viscous energy dissipation.
So $\Hone$ is the space of fields with finite energy and finite dissipation
rate, and $\Honezero$ imposes in addition the no slip condition at the walls.

\subsection{The three spaces of the mixed problem}

The mixed formulation has three unknowns and therefore three spaces.

\begin{definition}[Velocity space]
\[
  \Vsp = [\Honezero]^{2}
  = \left\{ \bv = (v_{1},v_{2}) : v_{1}, v_{2} \in \Honezero \right\} ,
  \qquad
  \normH{\bv}^{2} = \normH{v_{1}}^{2} + \normH{v_{2}}^{2} .
\]
\end{definition}

Note carefully what has changed from Lesson 2. There we used $\Vdiv$, which was
$\Vsp$ together with the extra requirement $\dive \bv = 0$. Now that requirement
has been removed from the space. The velocity space is the plain product of two
copies of $\Honezero$, with no constraint at all. The constraint has not been
abandoned, it has been moved: it will reappear in Section \ref{sec:mixed} as an
equation rather than as a property of the space.

\begin{definition}[Microrotation space]
\[
  \Wsp = \Honezero .
\]
\end{definition}

This is unchanged from Lesson 2. The microrotation is a scalar in two
dimensions, it satisfies a homogeneous Dirichlet condition, and its equation is
of reaction diffusion type, so $\Honezero$ is the natural space for exactly the
reasons given in Lesson 1.

\begin{definition}[Pressure space]\label{def:Q}
\[
  \Qsp = \Ltwozero = \left\{ q \in \Ltwo : \int_\Om q \dx = 0 \right\} ,
  \qquad
  \lVert q \rVert_{\Qsp} = \normL{q} .
\]
\end{definition}

This space is new, since in Lesson 2 there was no pressure unknown at all, so I
want to say carefully both why the pressure lives in $\Ltwo$ rather than
$\Hone$, and why we take the zero mean subspace.

\paragraph{Why $L^{2}$ and not $H^{1}$.}
The rule from Lesson 1 is that the space is read off from the weak formulation
by asking what regularity makes its integrals finite. In the mixed formulation
the pressure appears only in the term
\[
  - \int_\Om p \,\dive \bv \dx .
\]
There is no derivative on $p$ here. It was removed by integration by parts and
transferred onto $\bv$, which is exactly the trade we always make. For this
integral to be finite we need only $p \in \Ltwo$, since $\dive \bv \in \Ltwo$
whenever $\bv \in \Vsp$ and Cauchy--Schwarz then applies. Asking for $p \in
\Hone$ would be demanding more regularity than the formulation requires, and by
the argument of Lesson 1 that is the wrong thing to do. So the pressure sits in
a space with no derivatives at all, one step rougher than the velocity. This
asymmetry between the velocity space and the pressure space is characteristic of
mixed problems.

\paragraph{Why the zero mean condition.}
In the momentum equation the pressure appears only through $\grad p$. Hence if
$p$ solves the problem then so does $p + C$ for any constant $C$, because
$\grad(p+C) = \grad p$. The physics determines pressure differences, not
absolute pressure, so without some normalisation the solution is not unique.

The set of constants is a one dimensional subspace of $\Ltwo$, and $\Ltwo$
splits orthogonally as
\[
  \Ltwo = \Ltwozero \oplus \left\{ \text{constants} \right\} .
\]
To see this, given $q \in \Ltwo$ write
\[
  q = \underbrace{\left( q - \frac{1}{\abs{\Om}}\int_\Om q \dx \right)}_{\text{zero mean}}
  + \underbrace{\frac{1}{\abs{\Om}}\int_\Om q \dx}_{\text{constant}} ,
\]
and check that the first bracket integrates to zero while the two pieces are
orthogonal in $\Ltwo$. Restricting to $\Ltwozero$ therefore amounts to choosing
one representative from each family $\{p + C : C \in \R\}$, namely the one with
zero average.

\begin{remark}[What quotienting means, and where you have already seen it]
\label{rem:quotient}
To quotient a space by a subspace means to declare two elements equal whenever
they differ by an element of that subspace, and then to work with the resulting
families rather than with individual elements. Clock arithmetic is the familiar
example: $14$ and $2$ are the same time because they differ by a multiple of
$12$.

Here the subspace is the constants, and $p$ and $p + C$ are the same physical
state. Rather than carrying families around we pick one representative from each,
which is what Definition \ref{def:Q} does.

This is not an abstraction with no consequences. If the pressure space is left
as all of $\Ltwo$, so that the constants are still present, then the
corresponding matrix in a finite element code has a nontrivial null space: the
vector of nodal values representing a constant pressure is mapped to zero. The
matrix is then singular, and a direct solver will either fail or return
nonsense. The step of fixing the constant in code, whether by a Lagrange
multiplier or by pinning one pressure degree of freedom, is the discrete
counterpart of passing to $\Ltwozero$. So the abstract quotient and the singular
matrix are the same phenomenon seen from two sides.
\end{remark}

\subsection{The product spaces}

Following the packaging idea of Lesson 2, we group the unknowns. Write
\[
  \Vsp \times \Wsp
  \quad \text{with} \quad
  \lVert (\bv, z) \rVert_{\Vsp \times \Wsp}^{2} = \normH{\bv}^{2} + \normH{z}^{2} ,
\]
for the pair carrying the coercive part of the problem, and
\[
  \Xsp = (\Vsp \times \Wsp) \times \Qsp
  \quad \text{with} \quad
  \lVert (\bv, z, q) \rVert_{\Xsp}^{2}
  = \normH{\bv}^{2} + \normH{z}^{2} + \normL{q}^{2} ,
\]
for all three unknowns together. Each is a Hilbert space, being a finite product
of Hilbert spaces with the sum of squares norm.

I will use $\Vsp \times \Wsp$ throughout, and $\Xsp$ only in Section
\ref{sec:notcoercive}, where the point is precisely that working on $\Xsp$ does
not succeed.

%==============================================================================
\section{The mixed formulation}
\label{sec:mixed}
%==============================================================================

\subsection{The strong form}

The system, as in Lesson 2, is
\begin{align}
  -(\nu + \nu_{r}) \lap \bu + \grad p
    &= 2\nu_{r} \grad^{\perp} w + \bff , \label{eq:mom} \\
  \dive \bu &= 0 , \label{eq:incomp} \\
  -(c_{a} + c_{d}) \lap w + 4\nu_{r} w
    &= 2\nu_{r} \curl \bu + g , \label{eq:micro}
\end{align}
with $\bu = \bzero$ and $w = 0$ on $\dOm$.

\subsection{Task 3.A: deriving the mixed weak form}

The derivation follows Lesson 2 exactly, with one change: we do not discard the
pressure term, because the test functions are no longer divergence free.

\paragraph{Step 1: the momentum equation.}
Multiply \eqref{eq:mom} by $\bv \in \Vsp$ and integrate over $\Om$:
\[
  -(\nu+\nu_{r}) \int_\Om (\lap \bu)\cdot \bv \dx
  + \int_\Om \grad p \cdot \bv \dx
  = 2\nu_{r} \int_\Om \grad^{\perp} w \cdot \bv \dx
  + \int_\Om \bff \cdot \bv \dx .
\]

For the Laplacian term, apply Green's identity componentwise. The boundary terms
vanish because $\bv \in [\Honezero]^{2}$ has zero trace:
\[
  -(\nu+\nu_{r}) \int_\Om (\lap \bu)\cdot \bv \dx
  = (\nu+\nu_{r}) \ip{\grad \bu}{\grad \bv} .
\]

For the pressure term, apply the integration by parts formula
\[
  \int_\Om (\partial_{i}\phi)\psi \dx
  = \int_{\dOm} \phi \psi n_{i} \ds - \int_\Om \phi(\partial_{i}\psi) \dx
\]
with $\phi = p$ and $\psi = v_{i}$, once for each component. The boundary terms
vanish because $v_{i}$ has zero trace, so
\begin{equation}\label{eq:pressure-term}
  \int_\Om \grad p \cdot \bv \dx
  = \sum_{i=1}^{2} \int_\Om (\partial_{i}p) v_{i} \dx
  = - \sum_{i=1}^{2} \int_\Om p (\partial_{i}v_{i}) \dx
  = - \int_\Om p \,\dive \bv \dx .
\end{equation}

In Lesson 2 this was the end of the story, because $\dive \bv = 0$ made
\eqref{eq:pressure-term} vanish. Now $\bv$ is an arbitrary element of $\Vsp$, so
the term survives. Define
\begin{equation}\label{eq:bdef}
  b(\bv, q) := - \int_\Om q \,\dive \bv \dx .
\end{equation}
The momentum equation in weak form is then
\begin{equation}\label{eq:weakmom}
  (\nu+\nu_{r})\ip{\grad \bu}{\grad \bv} - 2\nu_{r}\ip{\grad^{\perp}w}{\bv}
  + b(\bv, p) = \ip{\bff}{\bv}
  \qquad \text{for all } \bv \in \Vsp .
\end{equation}

\paragraph{Step 2: the microrotation equation.}
Unchanged from Lesson 2. Multiply \eqref{eq:micro} by $z \in \Wsp$, integrate,
apply Green's identity to the Laplacian term with the boundary term vanishing
because $z$ has zero trace, and leave the zeroth order term alone since it
contains no derivative:
\begin{equation}\label{eq:weakmicro}
  (c_{a}+c_{d})\ip{\grad w}{\grad z} + 4\nu_{r}\ip{w}{z}
  - 2\nu_{r}\ip{\curl \bu}{z} = \ip{g}{z}
  \qquad \text{for all } z \in \Wsp .
\end{equation}

\paragraph{Step 3: the incompressibility equation.}
This is the genuinely new step. In Lesson 2 the constraint was built into the
space and never tested. Now we test it. Multiply \eqref{eq:incomp} by
$q \in \Qsp$ and integrate:
\[
  \int_\Om q \,\dive \bu \dx = 0
  \qquad \text{for all } q \in \Qsp ,
\]
which in the notation \eqref{eq:bdef} is
\begin{equation}\label{eq:weakincomp}
  b(\bu, q) = 0 \qquad \text{for all } q \in \Qsp .
\end{equation}
No integration by parts is needed, since there is nothing to move.

\paragraph{Step 4: adding the first two equations.}
As in Lesson 2, testing \eqref{eq:weakmom} against every $\bv$ and
\eqref{eq:weakmicro} against every $z$ is equivalent to testing their sum
against every pair $(\bv,z)$: choosing $z = 0$ recovers the first, choosing
$\bv = \bzero$ recovers the second. Writing $a$ for the bilinear form of Lesson
2,
\begin{multline}\label{eq:adef}
  a\big( (\bu,w),(\bv,z) \big)
  = (\nu+\nu_{r})\ip{\grad \bu}{\grad \bv}
  + (c_{a}+c_{d})\ip{\grad w}{\grad z}
  + 4\nu_{r}\ip{w}{z} \\
  {} - 2\nu_{r}\ip{\grad^{\perp}w}{\bv} - 2\nu_{r}\ip{\curl \bu}{z} ,
\end{multline}
and $F(\bv,z) = \ip{\bff}{\bv} + \ip{g}{z}$, we arrive at the mixed problem.

\begin{definition}[Mixed weak formulation]\label{def:mixed}
Find $\big( (\bu,w), p \big) \in (\Vsp \times \Wsp) \times \Qsp$ such that
\begin{align}
  a\big( (\bu,w),(\bv,z) \big) + b(\bv, p) &= F(\bv,z)
    &&\text{for all } (\bv,z) \in \Vsp \times \Wsp , \label{eq:mixed1} \\
  b(\bu, q) &= 0
    &&\text{for all } q \in \Qsp . \label{eq:mixed2}
\end{align}
\end{definition}

Note the structure. The form $a$ is exactly the one analysed in Lesson 2, with
no change whatsoever; the only difference is that it is now posed on the larger
space $\Vsp \times \Wsp$ rather than $\Vdiv \times \Wsp$. The new object is $b$,
which appears twice, once carrying the pressure into the momentum equation and
once expressing the constraint. A system with this shape, a coercive block plus
a constraint block appearing symmetrically, is called a saddle point problem.

%==============================================================================
\section{Task 3.B: the form is not coercive on $\Xsp$}
\label{sec:notcoercive}
%==============================================================================

The natural first attempt is to package all three unknowns into one bilinear
form on $\Xsp$ and apply Lax--Milgram as before. This section shows that it does
not work, and identifies the reason.

Define, for $(\bu,w,p)$ and $(\bv,z,q)$ in $\Xsp$,
\begin{equation}\label{eq:calA}
  \mathcal{A}\big( (\bu,w,p),(\bv,z,q) \big)
  := a\big( (\bu,w),(\bv,z) \big) + b(\bv,p) - b(\bu,q) .
\end{equation}
The minus sign on the last term is a convention, chosen so that the second
equation of Definition \ref{def:mixed} is recovered with the correct sign when
we test with $q$; any consistent choice would do.

\begin{proposition}[No coercivity on $\Xsp$]\label{prop:nocoerc}
There is no constant $\alpha > 0$ with
\[
  \mathcal{A}\big( (\bu,w,p),(\bu,w,p) \big)
  \ge \alpha \lVert (\bu,w,p) \rVert_{\Xsp}^{2}
  \qquad \text{for all } (\bu,w,p) \in \Xsp .
\]
\end{proposition}

\begin{proof}
First compute what happens on the diagonal. Setting $\bv = \bu$, $z = w$ and
$q = p$ in \eqref{eq:calA},
\[
  \mathcal{A}\big( (\bu,w,p),(\bu,w,p) \big)
  = a\big( (\bu,w),(\bu,w) \big) + b(\bu,p) - b(\bu,p)
  = a\big( (\bu,w),(\bu,w) \big) .
\]
The two occurrences of $b$ cancel exactly. This is not an accident of the sign
convention: whatever sign is chosen, the pressure enters $\mathcal{A}$ only
through $b$, and on the diagonal both occurrences have the same arguments, so
they cancel or double but in neither case produce anything involving
$\normL{p}^{2}$.

Now take the explicit choice
\[
  \bu = \bzero , \qquad w = 0 , \qquad p = q_{0}
\]
where $q_{0} \in \Qsp$ is any nonzero element, for instance
$q_{0}(x,y) = \cos(\pi x)$ on $\Om = (0,1)^{2}$, which has zero mean since
$\int_{0}^{1}\cos(\pi x) \dx = 0$ and is not identically zero.

For this choice the left hand side is
\[
  \mathcal{A}\big( (\bzero,0,q_{0}),(\bzero,0,q_{0}) \big)
  = a\big( (\bzero,0),(\bzero,0) \big) = 0 ,
\]
since $a$ is bilinear and therefore vanishes when either argument is zero. The
right hand side is
\[
  \alpha \lVert (\bzero,0,q_{0}) \rVert_{\Xsp}^{2}
  = \alpha \left( 0 + 0 + \normL{q_{0}}^{2} \right)
  = \alpha \normL{q_{0}}^{2} > 0 ,
\]
because $q_{0} \neq 0$. So the asserted inequality would read $0 \ge \alpha
\normL{q_{0}}^{2} > 0$, which is false for every $\alpha > 0$.
\end{proof}

\begin{remark}[Why this is structural rather than a failure of technique]
It is worth being clear that Proposition \ref{prop:nocoerc} is not saying that
we were not clever enough. Coercivity on $\Xsp$ would require the form to
control $\normL{p}^{2}$ from below, and the form simply contains no term in
which the pressure is paired with itself. Compare with the velocity, which
appears in $(\nu+\nu_{r})\ip{\grad \bu}{\grad \bv}$ and so produces
$\normL{\grad \bu}^{2}$ on the diagonal, or the microrotation, which appears in
$4\nu_{r}\ip{w}{z}$ and produces $\normL{w}^{2}$.

The reason is physical. Velocity and microrotation are diffusing fields, and the
diffusion terms represent dissipation of their energy. The pressure is not a
field of that kind at all. It is a Lagrange multiplier: a quantity whose only
role is to enforce the constraint $\dive \bu = 0$, adjusting itself to whatever
value is needed. It has no equation of its own, no diffusion, and no energy. A
Lagrange multiplier can never appear coercively, in any problem, so no amount of
manipulation will produce a $\normL{p}^{2}$ term.

The conclusion is that Lax--Milgram is simply the wrong theorem here, and we
need one designed for problems with constraints.
\end{remark}

%==============================================================================
\section{Brezzi's theorem}
%==============================================================================

\subsection{The kernel}

The replacement theorem is stated in terms of the set of velocities that satisfy
the constraint.

\begin{definition}[Kernel of $b$]\label{def:kernel}
\[
  \Zsp = \left\{ (\bv,z) \in \Vsp \times \Wsp
  : b(\bv,q) = 0 \ \text{ for all } q \in \Qsp \right\} .
\]
\end{definition}

\subsection{Statement}

\begin{theorem}[Brezzi]\label{thm:brezzi}
Let $\Vsp \times \Wsp$ and $\Qsp$ be Hilbert spaces, let $a$ be a bounded
bilinear form on $(\Vsp\times\Wsp)^{2}$ and $b$ a bounded bilinear form on
$(\Vsp\times\Wsp)\times\Qsp$. Suppose:

\begin{itemize}
\item[\textbf{(B1)}] \emph{Coercivity on the kernel.} There is $\alpha > 0$ with
\[
  a\big( (\bv,z),(\bv,z) \big)
  \ge \alpha \lVert (\bv,z) \rVert_{\Vsp\times\Wsp}^{2}
  \qquad \text{for all } (\bv,z) \in \Zsp .
\]

\item[\textbf{(B2)}] \emph{The inf-sup condition.} There is $\beta > 0$ with
\[
  \inf_{q \in \Qsp \setminus \{0\}} \
  \sup_{\bv \in \Vsp \setminus \{\bzero\}} \
  \frac{b(\bv,q)}{\normH{\bv}\,\normL{q}} \ \ge\ \beta .
\]
\end{itemize}

Then for every bounded linear functional $F$ the mixed problem of Definition
\ref{def:mixed} has exactly one solution $\big( (\bu,w),p \big)$, and it depends
continuously on the data, with constants depending only on $\alpha$, $\beta$ and
the boundedness constants of $a$ and $b$.
\end{theorem}

I am quoting this theorem rather than proving it. The proof is not
short and it belongs to the general theory of saddle point problems; Boffi,
Brezzi and Fortin give it in full.

\subsection{What has changed, and what has not}

Two hypotheses now replace the single coercivity hypothesis of Lax--Milgram.

Condition (B1) is weaker than full coercivity, because it is only required on
the kernel $\Zsp$ rather than on all of $\Vsp \times \Wsp$. That is exactly what
makes the theorem usable: we saw in Proposition \ref{prop:nocoerc} that
coercivity on the whole space is hopeless, but on the kernel the constraint is
satisfied and the situation is the one from Lesson 2.

Condition (B2) is entirely new and is the price of enlarging the space. It is
what replaces the information we lost when we stopped building the constraint
into the velocity space.

The important practical point is that Lesson 2 has already done the work for
(B1). Task 3.C below shows that $\Zsp = \Vdiv \times \Wsp$, which is precisely
the space on which Lesson 2 established coercivity with
\[
  \alpha = \frac{\min(\nu,\ c_{a}+c_{d})}{C_\Om^{2}+1} .
\]
So (B1) is already proved, and the whole of the remaining difficulty is (B2).

%==============================================================================
\section{Reading the inf-sup condition}
%==============================================================================

The inf-sup condition is the least transparent object in this subject on first
meeting, so this section takes it apart term by term.

\subsection{Unpacking the notation}

The expression is
\[
  \inf_{q \neq 0} \ \sup_{\bv \neq \bzero} \
  \frac{b(\bv,q)}{\normH{\bv}\,\normL{q}} \ \ge\ \beta .
\]

Read from the inside out. For a fixed $q$, the inner supremum asks: over all
possible velocity test functions $\bv$, how large can the quantity
$b(\bv,q)$ be made, once we normalise by the sizes of $\bv$ and $q$ so that the
answer does not depend on scaling? Call this number $S(q)$.

The outer infimum then asks: as $q$ ranges over all nonzero pressures, how small
can $S(q)$ get? The condition demands that it never drops below the fixed
positive number $\beta$.

Equivalently, and this form is usually easier to work with: for every nonzero
$q \in \Qsp$ there exists $\bv \in \Vsp$ with
\begin{equation}\label{eq:infsup-equiv}
  b(\bv,q) \ \ge\ \beta \normH{\bv}\,\normL{q} .
\end{equation}

\subsection{What it says in words}

Since $b(\bv,q) = -\int_\Om q \,\dive \bv \dx$, the condition says that
\emph{every pressure must be visible to some velocity}. For each candidate
pressure $q$, there has to be a velocity field whose divergence pairs with $q$
in a nontrivial way.

Suppose instead that some nonzero $q_{0}$ had $b(\bv,q_{0}) = 0$ for every
$\bv \in \Vsp$. Then $q_{0}$ would not appear anywhere in equation
\eqref{eq:mixed1}, since that equation only ever sees the pressure through
$b(\bv,\cdot)$. Consequently, if $\big((\bu,w),p\big)$ were a solution, so would
$\big((\bu,w),p+q_{0}\big)$ be, and uniqueness would fail. The inf-sup condition
is what rules this out.

\subsection{Why a uniform constant is needed}

It would be a weaker requirement to ask merely that for each nonzero $q$ there
exists some $\bv$ with $b(\bv,q) \neq 0$. That is a statement about each $q$
separately. The inf-sup condition asks for more: a single constant $\beta$ that
works for all $q$ simultaneously. This is what the word \emph{uniform} means
here.

The reason the stronger version is needed is as follows. Suppose the weaker
statement held but no uniform $\beta$ existed. Then there would be a sequence
$q_{n}$ of pressures, each visible to some velocity, but with
\[
  S(q_{n}) \longrightarrow 0 .
\]
Each $q_{n}$ is detected by the equations, but progressively more feebly. Now,
Brezzi's theorem produces a bound on the solution of the form
\[
  \normL{p} \le \frac{C}{\beta}\left( \text{data} \right) ,
\]
so as $\beta$ degenerates this bound blows up. The map taking the data $F$ to
the solution $p$ is called the \emph{solution operator}, and to say it is
\emph{unbounded} is to say that no such finite constant exists: arbitrarily
small data could produce arbitrarily large pressure. That is a failure of
stability, and a problem with an unbounded solution operator is not well posed
in any useful sense, even if a solution technically exists and is unique.

This has a very concrete numerical consequence, which is the subject of Lesson
4. In the discrete setting the constant $\beta_{h}$ for the finite element
spaces may degenerate as the mesh is refined, and when it does the computed
pressure develops large spurious oscillations even though the velocity may look
perfectly reasonable.

\subsection{Why the pressure space had to be $\Ltwozero$}

Task 3.E below shows that the inf-sup condition fails outright if the pressure
space is taken to be all of $\Ltwo$, because constant pressures are invisible to
every velocity in $[\Honezero]^{2}$. Passing to $\Ltwozero$ removes exactly those
invisible elements, and by the discussion in Remark \ref{rem:quotient} this is
the same step as fixing the singular matrix in a computation. The abstract
requirement and the practical fix are two views of one fact.

\subsection{Status of the continuous inf-sup condition}

For the specific form $b(\bv,q) = -\int_\Om q \,\dive \bv \dx$ on
$[\Honezero]^{2} \times \Ltwozero$, condition (B2) does hold on any bounded
Lipschitz domain. It is equivalent to the statement that the divergence operator
has a bounded right inverse: given $q \in \Ltwozero$, there exists
$\bv \in [\Honezero]^{2}$ with
\[
  \dive \bv = q \qquad \text{and} \qquad \normH{\bv} \le C \normL{q} .
\]
Constructing such a $\bv$ is a substantial theorem, associated with the names of
Bogovskii and of Ne\v{c}as. I am citing it, not proving it. Deciding what to
prove and what to quote is itself part of doing this properly, and this one is
firmly on the quoting side.

%==============================================================================
\section{Task 3.C: identifying the kernel}
%==============================================================================

\begin{proposition}\label{prop:kernel}
With $\Zsp$ as in Definition \ref{def:kernel},
\[
  \Zsp = \Vdiv \times \Wsp .
\]
\end{proposition}

Before the proof, note where the difficulty lies. We want to conclude from
\[
  \int_\Om q \,\dive \bv \dx = 0 \quad \text{for all } q \in \Qsp
\]
that $\dive \bv = 0$. The obvious move is to choose $q = \dive \bv$, which would
give $\normL{\dive \bv}^{2} = 0$ immediately. But that choice is only legitimate
if $\dive \bv$ actually belongs to $\Qsp = \Ltwozero$, and membership of
$\Ltwozero$ requires zero mean, which is not automatic. So the first step of the
proof must be to check that $\dive \bv$ does have zero mean.

\begin{proof}
The microrotation component plays no role, since $z$ does not appear in $b$ at
all. So $(\bv,z) \in \Zsp$ if and only if $z \in \Wsp$ and $b(\bv,q) = 0$ for
all $q \in \Qsp$. It therefore suffices to show that the latter condition is
equivalent to $\dive \bv = 0$.

\paragraph{Step 1: $\dive \bv$ has zero mean.}
Let $\bv \in \Vsp = [\Honezero]^{2}$. By the divergence theorem,
\[
  \int_\Om \dive \bv \dx = \int_{\dOm} \bv \cdot \bn \ds .
\]
Since $\bv$ has zero trace on $\dOm$, the integrand on the right vanishes
identically, so
\begin{equation}\label{eq:divzeromean}
  \int_\Om \dive \bv \dx = 0 .
\end{equation}
Also $\dive \bv \in \Ltwo$, because $\bv \in [\Hone]^{2}$ means each
$\partial_{i}v_{i} \in \Ltwo$ and a sum of $\Ltwo$ functions is in $\Ltwo$.
Combining these two facts,
\[
  \dive \bv \in \Ltwozero = \Qsp .
\]
So $\dive \bv$ is an admissible choice of test function, which is what we needed
to establish.

\paragraph{Step 2: the forward implication.}
Suppose $b(\bv,q) = 0$ for all $q \in \Qsp$. By Step 1 we may take
$q = \dive \bv$, giving
\[
  0 = b(\bv, \dive \bv) = -\int_\Om (\dive \bv)(\dive \bv) \dx
  = - \normL{\dive \bv}^{2} .
\]
Hence $\normL{\dive \bv} = 0$, and since $\normL{\cdot}$ is a norm this forces
$\dive \bv = 0$ as an element of $\Ltwo$, that is, almost everywhere in $\Om$.
Therefore $\bv \in \Vdiv$.

\paragraph{Step 3: the reverse implication.}
Conversely, if $\bv \in \Vdiv$ then $\dive \bv = 0$, so for every $q \in \Qsp$
\[
  b(\bv,q) = -\int_\Om q \cdot 0 \dx = 0 ,
\]
and $(\bv,z) \in \Zsp$ for any $z \in \Wsp$.

Combining Steps 2 and 3 gives $\Zsp = \Vdiv \times \Wsp$.
\end{proof}

\begin{remark}
This proposition is what connects the two lessons. Condition (B1) of Brezzi's
theorem asks for coercivity of $a$ on $\Zsp$, and Proposition \ref{prop:kernel}
identifies $\Zsp$ as $\Vdiv \times \Wsp$, which is exactly the space on which
Lesson 2 proved coercivity. So (B1) holds with
$\alpha = \min(\nu, c_{a}+c_{d})/(C_\Om^{2}+1)$, and nothing from Lesson 2 is
wasted. The whole of the analysis done there transfers to the mixed setting
without a single change.
\end{remark}

%==============================================================================
\section{Task 3.D: boundedness of $b$}
%==============================================================================

Brezzi's theorem requires $b$ to be bounded, so we check it. The proof needs one
small inequality which is worth isolating.

\begin{lemma}\label{lem:divbound}
For $\bv \in [\Hone]^{2}$,
\[
  \normL{\dive \bv} \le \sqrt{2}\, \normL{\grad \bv} .
\]
\end{lemma}

\begin{proof}
Pointwise, $\dive \bv = \partial_{1}v_{1} + \partial_{2}v_{2}$ is a sum of two
numbers. Apply the Cauchy--Schwarz inequality in $\R^{2}$ to the vectors
$(1,1)$ and $(\partial_{1}v_{1}, \partial_{2}v_{2})$:
\[
  \abs{\partial_{1}v_{1} + \partial_{2}v_{2}}
  \le \sqrt{1^{2}+1^{2}}\,
  \sqrt{(\partial_{1}v_{1})^{2} + (\partial_{2}v_{2})^{2}}
  = \sqrt{2}\,\sqrt{(\partial_{1}v_{1})^{2} + (\partial_{2}v_{2})^{2}} .
\]
Squaring and integrating over $\Om$,
\[
  \normL{\dive \bv}^{2}
  \le 2 \int_\Om \left[ (\partial_{1}v_{1})^{2} + (\partial_{2}v_{2})^{2} \right] \dx
  \le 2 \int_\Om \left[ (\partial_{1}v_{1})^{2} + (\partial_{2}v_{1})^{2}
  + (\partial_{1}v_{2})^{2} + (\partial_{2}v_{2})^{2} \right] \dx ,
\]
the last step being the discard of two non-negative terms, or rather the
addition of them. The final integral is exactly $\normL{\grad \bv}^{2}$, so
$\normL{\dive \bv}^{2} \le 2\normL{\grad \bv}^{2}$, and taking square roots
gives the claim.
\end{proof}

\begin{remark}
It is worth noting why Lemma \ref{lem:curldiv} of Lesson 2, which stated
$\normL{\grad \bv}^{2} = \normL{\curl \bv}^{2} + \normL{\dive \bv}^{2}$ and
would give the sharper bound $\normL{\dive \bv} \le \normL{\grad \bv}$, is not
used here. That identity required $\bv \in [\Honezero]^{2}$, because its proof
used integration by parts twice and needed the boundary terms to vanish. Lemma
\ref{lem:divbound} makes no use of boundary conditions and holds on all of
$[\Hone]^{2}$, at the cost of the constant $\sqrt{2}$. Since we will eventually
want boundedness of $b$ in settings where the boundary conditions may differ,
the more robust version is the right one to record.
\end{remark}

\begin{proposition}[Boundedness of $b$]\label{prop:bbounded}
For all $\bv \in \Vsp$ and $q \in \Qsp$,
\[
  \abs{b(\bv,q)} \le \sqrt{2}\, \normH{\bv}\, \normL{q} .
\]
\end{proposition}

\begin{proof}
By the definition \eqref{eq:bdef} and Cauchy--Schwarz in $\Ltwo$,
\[
  \abs{b(\bv,q)} = \abs{\int_\Om q \,\dive \bv \dx}
  \le \normL{q}\,\normL{\dive \bv} .
\]
By Lemma \ref{lem:divbound}, $\normL{\dive \bv} \le \sqrt{2}\normL{\grad \bv}$.
Finally $\normL{\grad \bv} \le \normH{\bv}$, since
$\normH{\bv}^{2} = \normL{\bv}^{2} + \normL{\grad \bv}^{2}$ and
$\normL{\bv}^{2} \ge 0$, so discarding a non-negative term gives the inequality.
Chaining the three estimates,
\[
  \abs{b(\bv,q)} \le \normL{q}\,\sqrt{2}\,\normL{\grad \bv}
  \le \sqrt{2}\,\normH{\bv}\,\normL{q} . \qedhere
\]
\end{proof}

%==============================================================================
\section{Task 3.E: failure of inf-sup on $\Ltwo$}
%==============================================================================

\begin{proposition}\label{prop:infsupfails}
Let the pressure space be taken as all of $\Ltwo$ rather than $\Ltwozero$. Then
the inf-sup condition fails: there is no $\beta > 0$ with
\[
  \sup_{\bv \in \Vsp \setminus \{\bzero\}}
  \frac{b(\bv,q)}{\normH{\bv}\normL{q}} \ge \beta
  \qquad \text{for all } q \in \Ltwo \setminus \{0\} .
\]
\end{proposition}

\begin{proof}
Take $q \equiv 1$, which is a nonzero element of $\Ltwo$ with
$\normL{1} = \abs{\Om}^{1/2} > 0$. For any $\bv \in \Vsp = [\Honezero]^{2}$,
\[
  b(\bv, 1) = -\int_\Om 1 \cdot \dive \bv \dx = - \int_\Om \dive \bv \dx .
\]
By the divergence theorem,
\[
  \int_\Om \dive \bv \dx = \int_{\dOm} \bv \cdot \bn \ds = 0 ,
\]
since $\bv$ has zero trace on $\dOm$. Hence $b(\bv,1) = 0$ for every
$\bv \in \Vsp$, and therefore
\[
  \sup_{\bv \neq \bzero} \frac{b(\bv,1)}{\normH{\bv}\normL{1}} = 0 .
\]
Since $\normL{1} > 0$, no positive $\beta$ can bound this supremum from below.
\end{proof}

\begin{remark}
This is the same calculation as Step 1 of Proposition \ref{prop:kernel}, seen
from the other side. There it showed that $\dive \bv$ automatically has zero
mean, which was what made the kernel argument work. Here the same fact shows
that constant pressures are orthogonal to the range of the divergence operator,
and so are invisible to the formulation.

The two are aspects of one statement: the divergence operator maps
$[\Honezero]^{2}$ into $\Ltwozero$ and not onto anything larger. Restricting the
pressure space to $\Ltwozero$ matches the pressure space to the range of the
operator, which is exactly what the inf-sup condition requires, and it is also
what removes the singular mode from the discrete matrix as described in Remark
\ref{rem:quotient}.
\end{remark}

%==============================================================================
\section{Summary and what comes next}
%==============================================================================

\subsection{Where the mixed problem stands}

Collecting the results of this lesson against the hypotheses of Theorem
\ref{thm:brezzi}:

\begin{center}
\begin{tabular}{@{}lll@{}}
\toprule
Requirement & Status & Reference \\
\midrule
$a$ bounded on $\Vsp\times\Wsp$ & proved & Lesson 2, Theorem on boundedness \\
$b$ bounded & proved & Proposition \ref{prop:bbounded} \\
Kernel identified as $\Vdiv\times\Wsp$ & proved & Proposition \ref{prop:kernel} \\
(B1) coercivity on the kernel & proved & Lesson 2, coercivity theorem \\
(B2) inf-sup condition & quoted & Bogovskii, Ne\v{c}as \\
\bottomrule
\end{tabular}
\end{center}

So the continuous mixed problem is well posed, and every ingredient except the
inf-sup condition itself was established by our own work.

\subsection{The shape of the argument, in one paragraph}

Enlarging the velocity space from $\Vdiv$ to $\Vsp$ made the problem
discretisable but destroyed coercivity, because the pressure entered as a
Lagrange multiplier and multipliers never appear coercively. Brezzi's theorem
repairs this by weakening the coercivity requirement to hold only on the kernel
of the constraint, where the original Lesson 2 analysis applies verbatim, and
adding in exchange a new condition ensuring that no pressure is invisible to the
velocity space. The zero mean normalisation of the pressure space is not a
convention but a necessity, since without it the constants form exactly such an
invisible set.

\subsection{What comes next}

Everything in this lesson concerned the continuous problem. The next lesson
concerns the discrete one, and the central point is that condition (B2) does
\emph{not} pass automatically to finite element subspaces. Choosing
$\Vsp_{h} \subset \Vsp$ and $\Qsp_{h} \subset \Qsp$ gives a discrete inf-sup
constant $\beta_{h}$ which may be small, or may degenerate as the mesh is
refined, even though the continuous $\beta$ is perfectly good.

This is the reason that equal order elements such as $P_{1}/P_{1}$ fail for
Stokes, producing the characteristic checkerboard pressure oscillation, and the
reason the Taylor--Hood pair $P_{2}/P_{1}$ is used instead. That is Lesson 4.

\end{document}
